\documentclass[11pt]{article}

\usepackage{amsmath,amssymb}
\usepackage{xurl}
\usepackage[colorlinks=true,linkcolor=blue,citecolor=blue,urlcolor=blue]{hyperref}

% Shared commands for notes in docs/paper-gaps/.

\DeclareUnicodeCharacter{2080}{\ifmmode{}_{0}\else\textsubscript{0}\fi}
\DeclareUnicodeCharacter{2081}{\ifmmode{}_{1}\else\textsubscript{1}\fi}
\DeclareUnicodeCharacter{2082}{\ifmmode{}_{2}\else\textsubscript{2}\fi}
\DeclareUnicodeCharacter{2099}{\ifmmode{}_{n}\else\textsubscript{n}\fi}

\newcommand{\Lean}{\textsc{Lean}}
\ExplSyntaxOn
\NewDocumentCommand{\leanid}{m}{
  \ifmmode
    \textsf{\detokenize{#1}}
  \else
    \group_begin:
    \urlstyle{sf}
    \tl_set:Nn \l_tmpa_tl {#1}
    \tl_replace_all:Nnn \l_tmpa_tl {\_} {_}
    \exp_args:NV \path \l_tmpa_tl
    \group_end:
  \fi
}
\ExplSyntaxOff
\providecommand{\tr}{\operatorname{tr}}
\newcommand{\tnleanIssueBase}{https://github.com/LionSR/TNLean/issues/}
\newcommand{\tnleanPrBase}{https://github.com/LionSR/TNLean/pull/}
\newcommand{\tnleanDocBase}{https://sirui-lu.com/TNLean/docs/}
\newcommand{\ghissue}[1]{\href{\tnleanIssueBase#1}{GitHub issue \##1}}
\newcommand{\ghpr}[1]{\href{\tnleanPrBase#1}{PR \##1}}
\newcommand{\doclink}[1]{\href{\tnleanDocBase#1.html}{\path{#1}}}

% Dirac notation and the identity operator, as in blueprint/src/macros/common.tex.
% \providecommand keeps note-local definitions authoritative.
\usepackage{dsfont}
\providecommand{\ket}[1]{|#1\rangle}
\providecommand{\bra}[1]{\langle #1|}
\providecommand{\braket}[2]{\langle #1 | #2 \rangle}
\providecommand{\ketbra}[2]{|#1\rangle\!\langle #2|}
\providecommand{\Id}{\mathds{1}}
\providecommand{\one}{\mathds{1}}
\providecommand{\C}{\mathbb{C}}
\providecommand{\MN}[1]{M_{#1}(\C)}
\providecommand{\MD}{M_D(\C)}

% External referencing for paper sources.  The build helper writes sanitized
% auxiliary files under build/paper-aux/ and excludes duplicated source labels.
\usepackage{xr}

% Import a generated paper auxiliary file with a caller-supplied label prefix.
% The candidate paths support builds from docs/paper-gaps/, the repository root,
% and build/paper-aux/ itself.
\newcommand{\importpaperaux}[2]{%
  \IfFileExists{../../build/paper-aux/#2.aux}{%
    \externaldocument[#1]{../../build/paper-aux/#2}%
  }{%
    \IfFileExists{build/paper-aux/#2.aux}{%
      \externaldocument[#1]{build/paper-aux/#2}%
    }{%
      \IfFileExists{#2.aux}{%
        \externaldocument[#1]{#2}%
      }{}%
    }%
  }%
}

% General external reference macros
\newcommand{\paperref}[2]{\ref{#1-#2}}
\newcommand{\papereqref}[2]{(\ref{#1-#2})}

% CPSV16 source (1606.00608 / MPDO-22-12-17-2.tex) external document and shortcuts
\importpaperaux{cpsv16-}{1606.00608/MPDO-22-12-17-2-tnlean}
\newcommand{\cpsvref}[1]{\paperref{cpsv16}{#1}}
\newcommand{\cpsveqref}[1]{\papereqref{cpsv16}{#1}}

% Verdict marker: \gapnote{<kind>}{<status>}. Kind is the policy
% classification (clarification, local-correction, scope-restriction,
% unfaithful, false-source, open-gap); status is open, wip (elimination
% underway), resolved, or historical. Machine-read by the site index;
% prints a verdict line.
\newcommand{\gapnote}[2]{\par\noindent\textbf{Verdict:} #1 (#2).\par}


\title{Phase-Normalized Weights in the Normal-Canonical Implication to Irreducible Form}
\author{}
\date{\today}

\begin{document}
\maketitle

\section{The assertion}

The periodic-MPS paper states that normal tensors are precisely periodic
tensors of period \(1\), and that when all blocks in the irreducible
decomposition are normal the decomposition is the canonical form in the sense of
Ci15.\footnote{\cite[Section~II]{DeLasCuevas2017Irreducible}; the local source
passage is arXiv:1708.00029, lines 258--271.}
In the same paragraph the irreducible-form weights are fixed by the convention
\(\mu_j>0\).

The blueprint records the corresponding implication: a normal canonical block family
assembles to an irreducible-form tensor whose periods are all equal to \(1\).

\section{The formal statements}

The formal predicate for irreducible form records the paper's positive-weight
convention: \(\operatorname{Re}\mu_j>0\) and \(\operatorname{Im}\mu_j=0\).
The formal predicate for normal canonical block hypotheses, inherited from the
non-periodic canonical-form formulation, assumes only that the complex weights
are nonzero and ordered by modulus.

The paper introduces a minimal non-repeated basis of periodic tensors later in
Section~II. That basis condition is separate from the irreducible-form
definition considered here.

The formalization now proves the phase-normalized construction.  From a normal
canonical block family with nonzero complex weights \(\mu_j\), it replaces
\(\mu_j\) by \(|\mu_j|\) and \(A_j\) by \((\mu_j/|\mu_j|)A_j\).  The assembled
tensor has the same matrix-product-vector family, the irreducible-form weights
are positive real numbers, and all periods are \(1\).\footnote{Formal
declarations:
\leanid{MPSTensor.toIsIrreducibleFormOfPhaseNormalized} and
\leanid{MPSTensor.toIsIrreducibleFormOfPhaseNormalized_period_eq_one} in
\doclink{TNLean/MPS/Periodic/NormalCanonicalPeriodOne.lean}.}

\section{Conclusion}

This note is now a record of the resolved discrepancy.  The paper's
irreducible-form convention uses positive weights, and the formalization
supplies the phase-normalization step needed to pass from the local
normal-canonical predicate, whose weights are nonzero complex numbers, to that
positive-weight convention without changing the represented
matrix-product-vector family.

\bibliographystyle{alpha}
\bibliography{references}

\end{document}
